% ============================================================================
%   CHAPTER 5 — RESULTS AND ANALYSIS
%   Aivancity: 10-15 pages — MUST BE COMPACT
%   Strategy: main results in body, all details in appendices
%
%   REORGANIZED to integrate mode collapse and non-trivial case analysis
%   REVISED to integrate the decoding-artefact reinterpretation of collapse,
%   the signal/prior mechanism, the re-thresholding intervention, and the
%   TREC external replication.
% ============================================================================

\chapter{Results and Analysis}
\label{chap:results}

% ----------------------------------------------------------------------------
% ----------------------------------------------------------------------------
\section{Chapter Overview}
\label{sec:res-overview}

This chapter reports the empirical findings of the study. The presentation 
is organized to distinguish confirmatory results (tests of the four 
pre-registered hypotheses stated in Section~\ref{sec:met-hypotheses}) from 
exploratory findings (novel observations that emerged during the analysis 
and are labelled here as original findings F1 through F9).

Sections~\ref{sec:res-audit} through~\ref{sec:res-h2} present the 
pre-LLM audit and the two pre-registered hypotheses that receive the 
strongest empirical support: H1 (systematic miscalibration, supported 
across all twelve configurations) and H2 (RLHF degrades calibration, 
supported by four independent permutation tests after Bonferroni 
correction). H2 is the central pre-registered contribution of this 
work.

Sections~\ref{sec:res-prompt} through~\ref{sec:res-non-trivial} report 
the three most novel findings, none of which was anticipated in the 
pre-registration. The enriched-prompt paradox 
(Section~\ref{sec:res-prompt}) shows that expert prompt engineering 
degrades calibration in five of six models. Mode collapse 
(Section~\ref{sec:res-mode-collapse}) documents that five of the twelve 
configurations concentrate more than 97\% of their predictions on a 
single label, invalidating the standard interpretation of their 
apparently reasonable ECE values. A deeper analysis of the underlying 
log-probabilities then shows that this collapse is a property of the 
decision rule rather than of the model's knowledge: the discriminative 
signal survives collapse and can be recovered by a non-monotone decision 
rule (Section~\ref{sec:res-rethreshold}). Performance on truly 
discriminative cases (Section~\ref{sec:res-non-trivial}) reveals that 
both medical LLMs tested achieve near-zero accuracy on the 238 cases 
requiring a genuine clinical decision, despite their apparent superiority 
on global accuracy.

Sections~\ref{sec:res-asymmetry} through~\ref{sec:res-h3h4} report 
secondary findings on the inclusion/exclusion asymmetry, model-specific 
evasion strategies, and the two pre-registered hypotheses that receive 
partial or no support (H3 on ordinal validity, H4 on selective 
classification). Sections~\ref{sec:res-recalibration} 
through~\ref{sec:res-baselines} present the post-hoc solutions --- 
isotonic recalibration, re-thresholding, and conformal prediction --- 
and the composite utility score that provides a single rank-ordering of 
the twelve configurations. Section~\ref{sec:res-external} reports an 
external replication of the collapse and its mechanism on the TREC 
Clinical Trials corpus. Section~\ref{sec:res-summary} synthesizes the 
outcomes of the four hypotheses and the twelve original findings.

Every point estimate reported in this chapter is accompanied by a 95\% 
bootstrap confidence interval, and every hypothesis test is reported 
with its Bonferroni-corrected significance threshold. Detailed tables, 
per-configuration reliability diagrams, and the full mode collapse 
analysis are provided in Appendices~\ref{app:tables}, 
\ref{app:reliability}, and~\ref{app:mode-collapse}.
% ----------------------------------------------------------------------------
% ----------------------------------------------------------------------------
\section{Pre-LLM Audit and Expert Agreement}
\label{sec:res-audit}

Before evaluating the twelve LLM configurations, we conduct a pre-LLM 
audit of the TrialGPT-Criterion-Annotations dataset itself. This audit 
serves two purposes: to establish the quality of the ground truth against 
which the LLMs will be compared, and to expose a systematic annotation 
bias in GPT-4 that has not been previously documented.

\subsection{Expert-GPT-4 Agreement}
\label{subsec:res-audit-kappa}

The dataset provides two independent annotations per pair: a consensus 
label from three physicians (which we treat as the ground truth 
throughout this thesis) and a label produced by GPT-4. The weighted 
kappa~\cite{cohen1968} between the two annotations, computed with 
quadratic weights on the ordinal scale, is $\kappa_w = 0.819$ 
(95\% bootstrap CI: [0.774, 0.86]). By conventional criteria for 
inter-rater agreement, this value falls in the range of substantial to 
almost perfect agreement, confirming that GPT-4 reproduces the 
physician consensus with high fidelity at the aggregate level.

This aggregate agreement, however, masks a systematic bias visible only 
at the label level. Table~\ref{tab:res-gpt4-bias} reports, for each of 
the six labels, the ratio of GPT-4 usage rate to expert usage rate. A 
ratio of 1.0 indicates that GPT-4 and experts use the label with the 
same frequency; a ratio above 1.0 indicates over-use by GPT-4; a ratio 
below 1.0 indicates under-use.

\begin{table}[htbp]
    \centering
    \caption{Usage-rate ratios of the six labels (GPT-4 vs.\ physician 
        consensus). A ratio of 1.0 indicates balanced usage; ratios far 
        from 1.0 reveal systematic annotation bias.}
    \label{tab:res-gpt4-bias}
    \small
    \begin{tabular}{@{}lccc@{}}
        \toprule
        \textbf{Label} & \textbf{Expert freq.} & \textbf{GPT-4 freq.} & \textbf{Ratio} \\
        \midrule
        not applicable         & 5.4\%  & 11.7\% & 2.16 \\
        not enough information & 28.9\% & 28.6\% & 0.99 \\
        included               & 7.3\%  & 6.9\%  & 0.95 \\
        not included           & 5.3\%  & 5.0\%  & 0.94 \\
        excluded               & 5.4\%  & 5.5\%  & 1.02 \\
        not excluded           & 47.7\% & 42.3\% & 0.89 \\
        \bottomrule
    \end{tabular}
\end{table}

Five of the six labels are used by GPT-4 within 10\% of the physician 
frequency, in line with the substantial overall kappa. The exception is 
``not applicable'', which GPT-4 uses at 2.16 times the physician rate. 
This asymmetric bias --- systematic over-use of one specific evasion 
label rather than a general tendency to over-use uncertainty --- is not 
reported in the original TrialGPT paper~\cite{jin2024trialgpt}, where 
GPT-4 is used as an annotation tool without such an audit. The bias is 
substantively meaningful: ``not applicable'' functions as a legitimate 
escape route when a criterion genuinely does not apply to a patient's 
situation (e.g., a pregnancy exclusion criterion applied to a male 
patient), but its over-use by GPT-4 suggests that the model treats a 
broader class of ambiguous cases as ``not applicable'' when a decisive 
label would have been warranted.

\begin{finding}[F1: GPT-4's Selective Evasion Bias]
    GPT-4 over-uses the ``not applicable'' label at 2.16 times the 
    physician rate, while using ``not enough information'' at the same 
    rate as the experts (ratio 0.99). This asymmetric pattern reveals a 
    selective evasion strategy that is not detected by aggregate 
    agreement metrics.
\end{finding}

This finding is meaningful for the downstream analysis of open-source 
LLMs in Section~\ref{sec:res-evasion}, where we show that open-source 
models exhibit a different but equally systematic evasion pattern --- 
over-use of ``not enough information'' rather than ``not applicable''. 
The two evasion channels appear to be model-specific, likely reflecting 
distinct alignment strategies during post-training.

\subsection{Ground Truth Choice}
\label{subsec:res-audit-choice}

Given the systematic bias identified in F1, we use the physician 
consensus rather than GPT-4 annotations as the ground truth throughout 
this thesis. This choice has three consequences. First, GPT-4 is not a 
model under evaluation in the calibration audit itself: its labels serve 
only as the input for F1. Second, the accuracy figures reported in 
subsequent sections are computed against expert labels, which increases 
their clinical relevance. Third, the class imbalance documented in 
Section~\ref{subsec:met-data-distribution} refers to the expert 
distribution and is preserved unchanged for the calibration analysis.
% ----------------------------------------------------------------------------
% ----------------------------------------------------------------------------
\section{Calibration Audit: H1 Supported Across All Configurations}
\label{sec:res-h1}

Hypothesis~H1 predicts that the Expected Calibration Error of an LLM-based 
eligibility scoring pipeline exceeds the well-calibrated threshold of 
$0.05$~\cite{vancalster2019calibration}. We test this hypothesis across all 
twelve configurations of the two-by-two matrix defined in 
Section~\ref{sec:met-models}, using the metrics and bootstrap protocol 
described in Section~\ref{sec:met-metrics}.

\subsection{Overall Calibration Results}
\label{subsec:res-h1-overall}

Table~\ref{tab:res-calibration} reports the ECE, its 95\% bootstrap 
confidence interval, and the Brier score for each of the twelve 
configurations. All twelve values fall well above the calibration 
threshold, with lower bounds of the 95\% intervals also strictly greater 
than 0.05. Hypothesis~H1 is therefore supported at the level of every 
individual configuration, not merely in aggregate.

\begin{table}[htbp]
    \centering
    \caption{Calibration metrics across the twelve configurations. All 
        ECE values exceed the 0.05 threshold, with the lower bound of the 
        95\% bootstrap confidence interval also strictly above 0.05 in 
        every case. Configurations are grouped by model family.}
    \label{tab:res-calibration}
    \small
    \begin{tabular}{@{}lccc@{}}
        \toprule
        \textbf{Configuration} & \textbf{ECE} & \textbf{95\% CI} & \textbf{Brier} \\
        \midrule
        \multicolumn{4}{@{}l}{\textit{Mistral-7B family}} \\
        Mistral-Base A     & 0.133 & [0.104, 0.163] & 0.258 \\
        Mistral-Base B     & 0.189 & [0.161, 0.218] & 0.269 \\
        Mistral-Instruct A & 0.322 & [0.297, 0.350] & 0.311 \\
        Mistral-Instruct B & 0.253 & [0.229, 0.287] & 0.353 \\
        \midrule
        \multicolumn{4}{@{}l}{\textit{Llama-3.1-8B family}} \\
        Llama-Base A       & 0.124 & [0.100, 0.148] & 0.187 \\
        Llama-Base B       & 0.151 & [0.123, 0.177] & 0.225 \\
        Llama-Instruct A   & 0.510 & [0.483, 0.535] & 0.456 \\
        Llama-Instruct B   & 0.602 & [0.575, 0.628] & 0.558 \\
        \midrule
        \multicolumn{4}{@{}l}{\textit{Medical LLMs}} \\
        Meditron A         & 0.182 & [0.153, 0.212] & 0.263 \\
        Meditron B         & 0.231 & [0.204, 0.259] & 0.283 \\
        BioMistral A       & 0.155 & [0.126, 0.182] & 0.234 \\
        BioMistral B       & 0.194 & [0.166, 0.222] & 0.244 \\
        \bottomrule
    \end{tabular}
\end{table}

The ECE values span a wide range, from $0.124$ (Llama-Base~A) to $0.602$ 
(Llama-Instruct~B). Even the best-calibrated configuration is more than 
twice the clinical threshold, and the worst is more than twelve times 
above it. The Brier score follows a similar ordering, confirming that 
the miscalibration is not compensated by favorable sharpness properties.

\medskip
\noindent\textbf{H1: Systematic Miscalibration Confirmed.}
All twelve configurations exhibit an ECE strictly greater than 0.05, 
with the lower bounds of their 95\% bootstrap confidence intervals 
also strictly above the threshold. Hypothesis~H1 is supported across 
the entire experimental matrix.
\medskip

\subsection{Direction of the Miscalibration}
\label{subsec:res-h1-direction}

ECE is a signed-magnitude metric: it measures how far confidence deviates 
from accuracy without indicating the direction of the deviation. To 
characterize the miscalibration more precisely, Table~\ref{tab:res-overconf} 
reports the mean overconfidence for each configuration, defined as the 
difference between mean confidence and mean accuracy 
(Equation~\ref{eq:overconf}). A positive value indicates that the model 
overstates its certainty on average; a negative value indicates 
understatement.

\begin{table}[htbp]
    \centering
    \caption{Mean overconfidence per configuration. Positive values 
        indicate that the model's mean confidence exceeds its mean 
        accuracy (overconfident); negative values indicate the opposite 
        (underconfident).}
    \label{tab:res-overconf}
    \small
    \begin{tabular}{@{}lc@{}}
        \toprule
        \textbf{Configuration} & \textbf{Overconfidence} \\
        \midrule
        Mistral-Base A     & $-0.189$ \\
        Mistral-Base B     & $-0.132$ \\
        Mistral-Instruct A & $+0.248$ \\
        Mistral-Instruct B & $+0.184$ \\
        Llama-Base A       & $+0.128$ \\
        Llama-Base B       & $+0.021$ \\
        Llama-Instruct A   & $+0.487$ \\
        Llama-Instruct B   & $+0.598$ \\
        Meditron A         & $-0.036$ \\
        Meditron B         & $+0.098$ \\
        BioMistral A       & $+0.144$ \\
        BioMistral B       & $+0.187$ \\
        \bottomrule
    \end{tabular}
\end{table}

The direction of miscalibration reveals a striking pattern: the Base 
variants of Mistral are systematically underconfident ($-0.189$ and 
$-0.132$), while their Instruct counterparts are systematically 
overconfident ($+0.248$ and $+0.184$). A similar transition occurs in 
the Llama family: Llama-Base is nearly calibrated on Prompt~B 
($+0.021$) but Llama-Instruct is severely overconfident ($+0.487$ and 
$+0.598$). This pattern anticipates the analysis of Section~\ref{sec:res-h2}, 
where we show that this transition from under-confidence to 
over-confidence is a direct effect of RLHF post-training.

\subsection{Reliability Diagrams}
\label{subsec:res-h1-reliability}

Figure~\ref{fig:res-reliability} shows the reliability diagrams for the 
twelve configurations, organized by model family and prompt variant. The 
diagonal line represents perfect calibration; points below the diagonal 
indicate overconfidence, points above indicate underconfidence.

\begin{figure}[htbp]
    \centering
    \makebox[\textwidth][c]{%
        \includegraphics[width=1.15\textwidth]{reliability_compare_AvsB.png}%
    }
    \caption{Reliability diagrams of the twelve configurations, grouped 
    by prompt variant (Prompt~A on the left, Prompt~B on the right), 
    with each curve corresponding to one of the six models. Each curve 
    plots the observed accuracy (labelled \emph{Précision observée}) 
    as a function of predicted confidence (\emph{Confiance prédite}) 
    in ten equal-width bins; the dashed reference line labelled 
    \emph{Parfait} corresponds to perfect calibration. Deviations from 
    the diagonal $y = x$ indicate miscalibration; the direction of the 
    deviation is consistent with the mean overconfidence reported in 
    Table~\ref{tab:res-overconf}. Axis labels and the legend are kept 
    in French, reflecting the working language of the experimental 
    phase.}
    \label{fig:res-reliability}
\end{figure}

Two visual patterns emerge from the reliability diagrams. First, the 
Mistral-Base and Meditron curves lie systematically above the diagonal 
in the low-confidence regions, indicating that these models are more 
accurate than they claim. Second, the Llama-Instruct and Mistral-Instruct 
curves lie systematically below the diagonal in the high-confidence 
regions, indicating that these models are less accurate than they 
claim. Per-configuration reliability diagrams are provided in 
Appendix~\ref{app:reliability} for detailed inspection.

% ----------------------------------------------------------------------------
% ----------------------------------------------------------------------------
\section{The Effect of RLHF: H2 Supported}
\label{sec:res-h2}

Hypothesis~H2 predicts that within a model family, the RLHF-aligned 
Instruct variant exhibits worse calibration than its Base counterpart. 
This is the central pre-registered contribution of this thesis: an 
extension of the observation of Kadavath et~al.~\cite{kadavath2022know}, 
who documented a similar effect on general question-answering tasks, to 
the specific domain of clinical trial matching. The comparison is 
performed within each of the two generalist families (Mistral-7B and 
Llama-3.1-8B), for each of the two prompt variants, yielding four 
independent tests of the same underlying hypothesis.

\subsection{Paired Comparison of ECE}
\label{subsec:res-h2-paired}

Table~\ref{tab:res-h2-paired} reports the paired ECE values for the four 
Base--Instruct comparisons. Under Prompt~A, both families show a 
substantial increase in ECE when moving from Base to Instruct: from 
$0.133$ to $0.322$ for Mistral, and from $0.124$ to $0.510$ for Llama. 
Under Prompt~B, the increase is smaller for Mistral (from $0.189$ to 
$0.253$) but larger for Llama (from $0.151$ to $0.602$). The direction 
of the effect is consistent across all four comparisons.

\begin{table}[htbp]
    \centering
    \caption{Paired ECE comparison for the four Base--Instruct tests of 
        Hypothesis~H2. $\Delta\ece$ is the difference $\ece(\text{Instruct}) 
        - \ece(\text{Base})$; a positive value supports H2. Permutation 
        p-values are computed with 1{,}000 resamples on the paired data; 
        the Bonferroni-corrected significance threshold across the eight 
        planned comparisons is $\alpha_{\text{corr}} = 0.00625$.}
    \label{tab:res-h2-paired}
    \small
    \begin{tabular}{@{}lccc>{\bfseries}c@{}}
        \toprule
        \textbf{Comparison} & 
        \textbf{ECE Base} & 
        \textbf{ECE Instruct} & 
        \boldmath$\Delta\ece$ & 
        \textbf{p-value} \\
        \midrule
        Mistral / Prompt A   & 0.133 & 0.322 & $+0.189$ & $<$0.001 \\
        Mistral / Prompt B   & 0.189 & 0.253 & $+0.064$ & 0.002 \\
        Llama / Prompt A     & 0.124 & 0.510 & $+0.386$ & $<$0.001 \\
        Llama / Prompt B     & 0.151 & 0.602 & $+0.451$ & $<$0.001 \\
        \bottomrule
    \end{tabular}
\end{table}

The four p-values are all below the Bonferroni-corrected threshold of 
$\alpha_{\text{corr}} = 0.05 / 8 = 0.00625$, which accounts for the eight 
planned comparisons in the pre-registered analysis plan 
(Section~\ref{subsec:met-bonferroni}). Even under this stringent 
correction, all four tests reject the null hypothesis of equal 
calibration between Base and Instruct variants.

\begin{finding}[F2: RLHF Systematically Degrades Calibration]
    Across two independent model families (Mistral-7B and Llama-3.1-8B) 
    and two independent prompt designs, RLHF-aligned Instruct versions 
    exhibit significantly worse calibration than their Base counterparts. 
    All four permutation tests reject the null hypothesis at 
    Bonferroni-corrected significance ($p < 0.003$), with $\Delta\ece$ 
    ranging from $+0.064$ to $+0.451$. Hypothesis~H2 is supported.
\end{finding}

\subsection{Mechanism: From Underconfidence to Overconfidence}
\label{subsec:res-h2-mechanism}

The ECE metric, as noted in Section~\ref{subsec:res-h1-direction}, is a 
signed-magnitude quantity and does not indicate the direction of the 
miscalibration. To characterize the mechanism by which RLHF degrades 
calibration, we compare the mean overconfidence of paired Base and 
Instruct variants (Table~\ref{tab:res-h2-mechanism}).

\begin{table}[htbp]
    \centering
    \caption{Mean overconfidence (mean confidence minus mean accuracy) 
        for the four Base--Instruct pairs. Negative values indicate 
        underconfidence, positive values indicate overconfidence. The 
        transition from negative to positive across every pair reveals a 
        consistent effect of RLHF on the direction of the miscalibration.}
    \label{tab:res-h2-mechanism}
    \small
    \begin{tabular}{@{}lccc@{}}
        \toprule
        \textbf{Comparison} & 
        \textbf{Base} & 
        \textbf{Instruct} & 
        \boldmath$\Delta$ \\
        \midrule
        Mistral / Prompt A & $-0.189$ & $+0.248$ & $+0.437$ \\
        Mistral / Prompt B & $-0.132$ & $+0.184$ & $+0.316$ \\
        Llama / Prompt A   & $+0.128$ & $+0.487$ & $+0.359$ \\
        Llama / Prompt B   & $+0.021$ & $+0.598$ & $+0.577$ \\
        \bottomrule
    \end{tabular}
\end{table}

The pattern is unambiguous: in every one of the four pairs, RLHF shifts 
the mean overconfidence in the positive direction, by amounts ranging 
from $+0.316$ to $+0.577$. The transition is most striking for Mistral, 
where the Base variants are systematically underconfident and the 
Instruct variants systematically overconfident. In effect, RLHF converts 
model modesty into model over-assertion.

This finding has a natural mechanistic interpretation. RLHF training 
optimizes for outputs that human raters prefer, and prior work has shown 
that human raters systematically prefer confident and assertive 
answers~\cite{ouyang2022training}. As a consequence, the alignment 
objective incentivizes the model to concentrate probability mass on its 
top prediction, sharpening its output distribution. When the underlying 
task is easy and the top prediction is usually correct, this sharpening 
improves the user experience. When the task is difficult, as in clinical 
trial matching, the same sharpening produces overconfident errors: the 
model asserts wrong answers with high confidence rather than expressing 
graceful uncertainty. Section~\ref{subsec:res-collapse-mechanism} 
refines this account by showing that the sharpening acts through a 
\emph{response prior} whose amplitude, relative to the discriminative 
signal, determines whether the model merely becomes overconfident or 
collapses entirely.

\subsection{Comparison to Prior Work}
\label{subsec:res-h2-comparison}

Kadavath et~al.~\cite{kadavath2022know} first documented that RLHF 
degrades calibration on general question-answering benchmarks (TriviaQA, 
Lambada, TruthfulQA), reporting an ECE increase from 0.02 to approximately 
0.15 after RLHF post-training. Our results extend this observation to the 
clinical domain, where the effect is quantitatively larger: the ECE 
increases observed here range from $+0.064$ to $+0.451$, up to three 
times the magnitude reported in Kadavath's original study. Two factors 
plausibly contribute to this amplification. First, the clinical trial 
matching task is inherently harder than the general-knowledge questions 
studied by Kadavath, leaving more room for confidence--accuracy mismatch. 
Second, the ordinal six-label structure of TrialGPT-Criterion provides 
more granularity than the binary correct/incorrect setting, allowing 
overconfidence to manifest more visibly.

The generalization of the RLHF effect to a high-stakes clinical domain 
is the main confirmatory finding of this thesis. Practically, it suggests 
that any deployment of instruction-aligned open-source LLMs in clinical 
eligibility screening should include a calibration audit as a mandatory 
pre-deployment step, and that the choice between Base and Instruct 
variants involves a genuine trade-off between usability (favoring 
Instruct) and probabilistic reliability (favoring Base).
% ----------------------------------------------------------------------------
% ----------------------------------------------------------------------------
\section{The Enriched-Prompt Paradox}
\label{sec:res-prompt}

Prior to running the experiments, we expected the enriched Prompt~B to 
improve calibration across the board: the addition of label definitions, 
worked examples, and the exclusion-inference rule was designed precisely 
to reduce ambiguity and disambiguate cases where the models were 
predicted to struggle. The observed pattern contradicts this expectation.

\subsection{Effect of Prompt Enrichment Across Models}
\label{subsec:res-prompt-comparison}

Table~\ref{tab:res-prompt-comparison} reports the paired ECE values for 
each of the six models under Prompt~A and Prompt~B, together with the 
signed difference $\Delta\ece = \ece_B - \ece_A$. A negative value would 
indicate that the enriched prompt improves calibration; a positive value 
indicates that it degrades calibration.

\begin{table}[htbp]
    \centering
    \caption{Effect of prompt enrichment on ECE. A positive 
        $\Delta\ece$ indicates that Prompt~B, despite being enriched with 
        expert guidance, produces \emph{worse} calibration than the 
        minimal Prompt~A.}
    \label{tab:res-prompt-comparison}
    \small
    \begin{tabular}{@{}lccc@{}}
        \toprule
        \textbf{Model} & 
        \textbf{ECE Prompt A} & 
        \textbf{ECE Prompt B} & 
        \boldmath$\Delta\ece$ \\
        \midrule
        Mistral-Base       & 0.133 & 0.189 & $+0.056$ \\
        Mistral-Instruct   & 0.322 & 0.253 & $-0.069$ \\
        Llama-Base         & 0.124 & 0.151 & $+0.027$ \\
        Llama-Instruct     & 0.510 & 0.602 & $+0.092$ \\
        Meditron           & 0.182 & 0.231 & $+0.049$ \\
        BioMistral         & 0.155 & 0.194 & $+0.039$ \\
        \bottomrule
    \end{tabular}
\end{table}

Five of the six models show a positive $\Delta\ece$: the enriched prompt 
degrades their calibration. The single exception is Mistral-Instruct, 
which is the only model for which Prompt~B produces an improvement 
($\Delta\ece = -0.069$). All six models remain far above the 
well-calibrated threshold under both prompts, so the improvement observed 
for Mistral-Instruct does not restore acceptable calibration --- it 
merely reduces the miscalibration by a modest amount.

\begin{finding}[F3: The Enriched-Prompt Paradox]
    Contrary to the intuition that expert prompt design improves 
    calibration, prompts enriched with label definitions, examples, and 
    inference rules \emph{degrade} calibration in five of six models. 
    Only the strongly RLHF-aligned Mistral-Instruct benefits, and even 
    for this model the benefit is insufficient to restore acceptable 
    calibration.
\end{finding}

\subsection{Interpretation}
\label{subsec:res-prompt-interpretation}

Two interpretations account for this paradoxical pattern. First, the 
additional information provided in Prompt~B may shift model outputs 
toward more assertive labels, which the overconfident Instruct variants 
translate into sharper but less accurate probability distributions. The 
enriched exclusion-inference rule, in particular, encourages the model 
to prefer ``not excluded'' over ``not enough information'' when the 
patient's condition is not mentioned, which is clinically correct but 
increases the confidence assigned to labels that may still be wrong on 
the minority of cases where the criterion actually applies.

Second, the enrichment interacts differently with model classes. Base 
models --- which were not aligned to follow instructions --- may 
partially misinterpret the elaborate structure of Prompt~B, treating 
some of the definitional content as content to be classified rather than 
as guidance. Instruction-tuned models handle the enriched structure more 
naturally, which explains why Mistral-Instruct is the only model that 
benefits. The dominant effect, however, remains negative: enrichment 
degrades more models than it helps, contradicting the standard 
recommendation to invest in prompt engineering as a route to better 
model behavior.
% ----------------------------------------------------------------------------
% ----------------------------------------------------------------------------
\section{Mode Collapse in Medical and RLHF-Aligned LLMs}
\label{sec:res-mode-collapse}

The calibration metrics reported so far implicitly assume that each 
configuration produces a diverse distribution of predictions across the 
six eligibility labels. This assumption breaks down for five of the 
twelve configurations, which concentrate almost all of their predictions 
on a single label. We term this behavior \emph{mode collapse}, and 
document it here as one of the most consequential findings of the study. 
Sections~\ref{subsec:res-collapse-distribution} 
and~\ref{subsec:res-collapse-structural} establish the phenomenon and its 
persistence; Sections~\ref{subsec:res-collapse-decoding} 
and~\ref{subsec:res-collapse-mechanism} then show that it is a property 
of the decision rule rather than of the model's knowledge.

\subsection{Prediction Distribution Analysis}
\label{subsec:res-collapse-distribution}

Table~\ref{tab:res-collapse-dist} reports, for each configuration, the 
proportion of predictions assigned to the modal (most frequent) label and 
the identity of that label. A ratio close to $1.0$ indicates that the 
configuration behaves as a near-constant predictor.

\begin{table}[htbp]
    \centering
    \caption{Modal label and its usage rate for each configuration. 
        Configurations with a modal usage above 90\% no longer function 
        as classifiers; they behave as biased constant predictors.}
    \label{tab:res-collapse-dist}
    \small
    \begin{tabular}{@{}lcc@{}}
        \toprule
        \textbf{Configuration} & 
        \textbf{Modal label} & 
        \textbf{Usage rate} \\
        \midrule
        BioMistral B       & not enough information & 100.0\% \\
        Llama-Instruct B   & not enough information & \phantom{0}99.7\% \\
        BioMistral A       & not enough information & \phantom{0}98.2\% \\
        Llama-Instruct A   & not enough information & \phantom{0}97.8\% \\
        Llama-Base B       & not enough information & \phantom{0}97.2\% \\
        \midrule
        Meditron A         & not excluded           & \phantom{0}64.8\% \\
        Meditron B         & not excluded           & \phantom{0}59.3\% \\
        \midrule
        Mistral-Base B     & not excluded           & \phantom{0}45.1\% \\
        Mistral-Base A     & not excluded           & \phantom{0}47.2\% \\
        Llama-Base A       & excluded               & \phantom{0}45.4\% \\
        Mistral-Instruct B & excluded               & \phantom{0}32.5\% \\
        Mistral-Instruct A & not enough information & \phantom{0}47.4\% \\
        \bottomrule
    \end{tabular}
\end{table}

Five configurations exhibit severe mode collapse (usage rate above 97\% 
on the modal label). Two additional configurations exhibit partial 
collapse (Meditron A and B, with usage rates of 64.8\% and 59.3\%). The 
remaining five configurations produce more diverse predictions, though 
none of them exceeds 50\% of the majority class distribution.

\begin{finding}[F7: Mode Collapse in Five of Twelve Configurations]
    Five open-source configurations --- both BioMistral variants, all 
    three Llama configurations with Prompt~B or with instruction tuning, 
    and to a lesser extent Meditron --- predict a single label in more 
    than 97\% of cases. Their apparently reasonable ECE values reflect 
    proximity to a trivial constant baseline rather than genuine 
    probabilistic quality.
\end{finding}

\subsection{Structural Persistence Under Constrained Decoding}
\label{subsec:res-collapse-structural}

An important question is whether mode collapse reflects a genuine absence 
of clinical reasoning, or merely an artifact of the label vocabulary. To 
distinguish these two possibilities, we conducted a controlled 
simulation: for each collapsed configuration, we restricted the softmax 
normalization (Equation~\ref{eq:softmax}) to the two decisive labels of 
each criterion type ($\{$included, not included$\}$ for inclusion 
criteria, $\{$excluded, not excluded$\}$ for exclusion criteria), 
effectively removing the evasion labels from the candidate set. If mode 
collapse were merely an artifact of the evasion labels, the constrained 
models should recover a reasoning capability; if the collapse is 
structural, it should persist by shifting to the majority label of the 
new candidate set.

The results, reported in Appendix~\ref{app:mode-collapse}, are 
unambiguous. Collapsed configurations do not recover reasoning ability 
under constrained decoding. Instead, their predictions migrate to the 
majority label of the criterion type: BioMistral, previously predicting 
``not enough information'', shifts to predicting ``not excluded'' in 
approximately 74\% of cases when restricted to two labels. Llama-Instruct 
shifts similarly. The collapse thus persists under a change of candidate 
set, confirming that evasion labels are not the cause of the collapse. As 
the next two subsections show, however, this persistence does not mean 
the model lacks the underlying signal: the collapse is a property of the 
$\arg\max$ decision rule, not of the model's knowledge.

\begin{finding}[F8: Mode Collapse Is Structural, Not Superficial]
    When forced to choose between the two decisive labels of each 
    criterion type, collapsed configurations do not recover reasoning 
    ability. Their collapse migrates from ``not enough information'' to 
    the majority label of the restricted candidate set (``not excluded'' 
    for exclusion criteria), showing that the collapse is not an artifact 
    of the evasion vocabulary but a property of the model's decision under 
    $\arg\max$.
\end{finding}

\subsection{Mode Collapse Is a Decoding Artefact, Not a Knowledge Loss}
\label{subsec:res-collapse-decoding}

Findings F7 and F8 establish that collapsed configurations concentrate 
their predictions on a single label and do not recover under constrained 
decoding. A natural reading of this behaviour is that these models lack 
the clinical knowledge required to discriminate. A finer analysis of the 
underlying probability distributions contradicts that reading: the 
discriminative signal is present in the log-probabilities, but its 
amplitude is too small to move the $\arg\max$.

Two observations support this interpretation. First, the collapse occurs 
at \emph{high} predictive entropy. BioMistral~B concentrates 100\% of its 
predictions on a single label, yet its mean normalized entropy is $0.888$ 
and its mean maximum probability only $0.483$: the output distribution is 
nearly flat, but its ordering never changes across the 1{,}015 pairs. 
This is not the sharpened, overconfident distribution one associates with 
collapse, but a distribution whose $\arg\max$ is frozen while its mass 
remains diffuse.

Second, the models rank cases correctly even when they never vary their 
predicted label. For each configuration we computed Cliff's~$\delta$ 
between the probability assigned to the positive label 
($p_{\text{included}}$ for inclusion criteria) on gold-positive versus 
gold-negative cases. Llama-Instruct~A, which predicts a single label in 
97.8\% of cases, obtains $\delta = 0.823$ ($p < 10^{-39}$) --- the 
highest ordinal separation of the entire study. The model knows which 
patients satisfy the criterion; it simply never expresses this knowledge 
in its decision. Table~\ref{tab:res-signal-expression} reports the 
ordinal separation for the collapsed and partially-collapsed 
configurations.

\begin{table}[htbp]
    \centering
    \caption{Ordinal separation (Cliff's $\delta$ on the positive-label 
        probability, inclusion criteria) for the collapsed configurations. 
        A high $\delta$ combined with a near-degenerate predicted-label 
        distribution indicates that the discriminative signal is present 
        but not expressed in the $\arg\max$. Meditron is the sole 
        configuration with a negative $\delta$.}
    \label{tab:res-signal-expression}
    \small
    \begin{tabular}{@{}lccc@{}}
        \toprule
        \textbf{Configuration} & 
        \textbf{Dominant rate} & 
        \textbf{Cliff's $\delta$} & 
        \textbf{$p$-value} \\
        \midrule
        Llama-Instruct A & 97.8\%  & $+0.823$ & $<10^{-39}$ \\
        Llama-Instruct B & 99.7\%  & $+0.733$ & $<10^{-31}$ \\
        BioMistral B     & 100.0\% & $+0.499$ & $<10^{-15}$ \\
        Llama-Base B     & 97.2\%  & $+0.353$ & $<10^{-7}$  \\
        BioMistral A     & 98.2\%  & $+0.103$ & $0.098$     \\
        Meditron A       & 64.8\%  & $-0.138$ & $0.026$     \\
        \bottomrule
    \end{tabular}
\end{table}

The distinction is decisive for the remainder of this chapter. Four of 
the collapsed configurations (both Llama-Instruct variants, BioMistral~B, 
Llama-Base~B) carry a substantial, statistically significant ordinal 
signal that the $\arg\max$ discards. Two configurations (Meditron~A 
and~B) carry no exploitable signal --- Meditron~A's separation is 
outright negative ($\delta = -0.138$), indicating that it ranks cases in 
the wrong order, and Meditron~B's is only marginally positive, far below 
that of any recoverable configuration. We return to this separation in 
Section~\ref{subsec:res-rethreshold-diag} and show that it maps onto two 
mechanistically distinct pathologies: recoverable collapse and signal 
corruption.

\subsection{A Mechanism for Collapse: Response Prior Versus Signal 
    Amplitude}
\label{subsec:res-collapse-mechanism}

If the discriminative signal survives collapse, why does the $\arg\max$ 
not reflect it? The answer lies in the relative magnitude of two 
quantities in the log-probability space: the amplitude with which the 
model varies the log-probability of the positive label across inputs 
(the \emph{signal}), and the systematic offset by which competing labels 
are preferred (the \emph{response prior}).

For each configuration we measured the signal as the standard deviation 
of the positive-label log-probability across pairs 
($\sigma_{\text{pos}}$), and the prior as the mean gap between the 
strongest competing label and the positive label, computed 
\emph{excluding} the positive label itself to avoid circularity 
($\text{bias}_{\text{ext}}$). Their ratio, 
$r = \sigma_{\text{pos}} / |\text{bias}_{\text{ext}}|$, orders the twelve 
configurations almost perfectly by their $\arg\max$ macro-F1 on inclusion 
criteria (Spearman $\rho = 0.866$, $p = 0.00027$). Every configuration 
with $r < 0.6$ collapses ($F_1^{\arg\max} = 0$); every configuration with 
$r > 0.75$ classifies ($F_1^{\arg\max} > 0.45$). No configuration falls 
between these two regimes.

\paragraph{Mechanism.} The ratio reorganizes the twelve configurations 
along a single axis. In the Llama family, RLHF drives the ratio down 
across the collapse threshold: Llama-Base~A sits in the healthy regime 
while Llama-Instruct~A falls to $r = 0.35$, deep in the collapsed regime 
(Table~\ref{tab:res-diagnostic}). In the Mistral family, the ratio stays 
above the threshold for every variant, base and instruct alike, and 
discrimination is preserved. The contrast is not that RLHF is harmful 
per se, but that it moves a configuration's ratio, and collapse follows 
whenever that movement pushes the response prior to dominate the signal. 
This refines Finding~F2: RLHF does not erase clinical knowledge, it 
reshapes the balance between signal and prior that the $\arg\max$ 
resolves, and when the prior wins, the decision saturates and masks the 
knowledge the model still holds. The per-configuration signal amplitude 
and response prior underlying each ratio are tabulated in 
Appendix~\ref{app:mode-collapse}.

This mechanism also explains, as a mathematical consequence rather than 
an empirical observation, why post-hoc recalibration cannot repair 
collapse (Section~\ref{subsec:res-recalibration-caveats}): any monotone 
recalibration preserves the $\arg\max$, hence leaves $r$ and the collapse 
untouched. It further predicts that a non-monotone decision rule --- one 
that shifts the decision boundary rather than rescaling confidence --- 
should recover the masked signal. Section~\ref{sec:res-rethreshold} tests 
this prediction directly.
% ----------------------------------------------------------------------------
% ----------------------------------------------------------------------------
\section{Performance on Truly Discriminative Cases}
\label{sec:res-non-trivial}

The class imbalance documented in Section~\ref{subsec:met-data-distribution} 
implies that a large fraction of the dataset can be classified correctly 
by trivial strategies. On the 47.7\% of pairs whose expert label is 
``not excluded'' and the 28.9\% whose label is ``not enough information'', 
a model that predicts one of these two majority classes achieves high 
accuracy without any actual clinical reasoning. The remaining 238 pairs 
(23.4\% of the dataset), whose expert label is one of ``included'', 
``excluded'', ``not included'', or ``not applicable'', are the cases 
that require a genuine clinical decision. We refer to these as the 
\emph{truly discriminative} cases.

\subsection{Comparison of Global vs.\ Discriminative Accuracy}
\label{subsec:res-non-trivial-comparison}

Table~\ref{tab:res-non-trivial} reports, for each configuration, the 
global accuracy across the 1{,}015 pairs and the accuracy on the 238 
truly discriminative cases. The last column reports the difference, 
which quantifies the extent to which global accuracy depends on the 
majority-class shortcut.

\begin{table}[htbp]
    \centering
    \caption{Global accuracy compared to accuracy on the 238 truly 
        discriminative cases (expert label outside 
        $\{$not excluded, not enough information$\}$). A large negative 
        drop indicates that the model's global accuracy depends on the 
        majority-class distribution rather than on clinical reasoning.}
    \label{tab:res-non-trivial}
    \small
    \begin{tabular}{@{}lccc@{}}
        \toprule
        \textbf{Configuration} & 
        \textbf{Global acc.} & 
        \textbf{Discriminative acc.} & 
        \textbf{Drop} \\
        \midrule
        Meditron A         & 0.648 & 0.000 & $-0.648$ \\
        Meditron B         & 0.630 & 0.000 & $-0.630$ \\
        BioMistral B       & 0.289 & 0.000 & $-0.289$ \\
        BioMistral A       & 0.296 & 0.004 & $-0.292$ \\
        Llama-Instruct B   & 0.292 & 0.013 & $-0.279$ \\
        Llama-Base B       & 0.308 & 0.013 & $-0.295$ \\
        Llama-Instruct A   & 0.290 & 0.021 & $-0.269$ \\
        Llama-Base A       & 0.230 & 0.265 & $+0.035$ \\
        Mistral-Base B     & 0.586 & 0.416 & $-0.170$ \\
        Mistral-Instruct B & 0.480 & 0.500 & $+0.020$ \\
        Mistral-Base A     & 0.500 & 0.613 & $+0.113$ \\
        \midrule
        \textbf{Mistral-Instruct A} & 0.308 & \textbf{0.719} & \textbf{$+0.411$} \\
        \bottomrule
    \end{tabular}
\end{table}

The results reverse the model hierarchy suggested by global accuracy. 
Meditron A, which appears to be the best-performing configuration in 
Table~\ref{tab:res-calibration} with a global accuracy of 64.8\%, drops 
to \textit{zero} accuracy on the 238 discriminative cases: it never 
predicts any of the four decisive labels correctly. BioMistral similarly 
achieves 0.0\% and 0.4\% on the discriminative subset. Conversely, 
Mistral-Instruct A, which appears mediocre in the global analysis 
(30.8\%), reaches 71.9\% accuracy on the cases that require actual 
clinical reasoning.

\begin{finding}[F9: Global Accuracy Hides Complete Failure on 
    Discriminative Cases]
    On the 238 pairs that require a truly discriminative clinical 
    decision, both medical LLMs tested (Meditron and BioMistral) achieve 
    accuracy at or below 0.4\%, despite global accuracy figures of 
    28--65\%. Only Mistral-Instruct with Prompt~A reaches clinically 
    acceptable performance on this subset (71.9\% accuracy). Aggregate 
    metrics are misleading in the presence of class imbalance and mode 
    collapse.
\end{finding}

\paragraph{Two distinct pathologies among the failing medical models.} 
Finding~F9 groups Meditron and BioMistral as jointly failing on the 
discriminative subset, but the two models fail for different reasons, a 
distinction developed in Section~\ref{subsec:res-collapse-decoding}. 
BioMistral retains a recoverable discriminative signal (area under the 
ROC curve $= 0.750$ on inclusion criteria, Cliff's $\delta > 0$): its 
zero discriminative accuracy is a decoding failure, not an absence of 
knowledge. Meditron~A, by contrast, has an area under the curve of 
$0.435$ --- below chance --- and a negative ordinal separation 
(Table~\ref{tab:res-signal-expression}): it ranks discriminative cases in 
the wrong order. Meditron thus attains the highest global accuracy of the 
study (64.8\%) while being its most degenerate configuration, winning on 
aggregate metrics purely by betting on the majority class. This 
separation --- recoverable collapse versus signal corruption --- is 
formalized as a two-dimensional diagnostic in 
Section~\ref{subsec:res-rethreshold-diag}.

\subsection{Implications for Model Evaluation}
\label{subsec:res-non-trivial-implications}

This finding has direct methodological implications for the evaluation 
of medical LLMs. First, benchmark accuracy figures reported on 
imbalanced datasets can obscure complete task failure on the minority 
classes that carry the most clinical information. Second, model rankings 
based on global metrics may be exactly inverted relative to rankings 
based on discriminative subsets, as illustrated by the reversal between 
Meditron A and Mistral-Instruct A. Third, the composite utility score 
introduced in Section~\ref{sec:res-baselines} incorporates prediction 
diversity as a factor precisely to prevent this class of masking effect. 
A responsible evaluation protocol for clinical LLMs should systematically 
report stratified accuracy alongside aggregate figures.
% ----------------------------------------------------------------------------
% ----------------------------------------------------------------------------
\section{Systematic Inclusion vs.\ Exclusion Asymmetry}
\label{sec:res-asymmetry}

The two criterion types --- inclusion and exclusion --- are structurally 
symmetric: both require a binary decision about whether the patient does 
or does not satisfy the criterion. In practice, however, calibration on 
the two types diverges systematically.

Table~\ref{tab:res-asymmetry} reports the ECE stratified by criterion 
type for each of the twelve configurations. The last column reports the 
ratio $\ece_{\text{excl}} / \ece_{\text{incl}}$, which quantifies the 
asymmetry.

\begin{table}[htbp]
    \centering
    \caption{ECE stratified by criterion type. Ratios above $1$ indicate 
        that calibration is worse on exclusion criteria than on inclusion 
        criteria. \model{Meditron} is the only model to reverse this 
        pattern.}
    \label{tab:res-asymmetry}
    \small
    \begin{tabular}{@{}lccc@{}}
        \toprule
        \textbf{Configuration} & 
        \textbf{ECE inclusion} & 
        \textbf{ECE exclusion} & 
        \textbf{Ratio} \\
        \midrule
        Mistral-Base A     & 0.089 & 0.198 & 2.22 \\
        Mistral-Base B     & 0.121 & 0.267 & 2.21 \\
        Mistral-Instruct A & 0.152 & 0.451 & 2.97 \\
        Mistral-Instruct B & 0.108 & 0.398 & 3.69 \\
        Llama-Base A       & 0.078 & 0.187 & 2.40 \\
        Llama-Base B       & 0.093 & 0.225 & 2.42 \\
        Llama-Instruct A   & 0.203 & 0.703 & 3.46 \\
        Llama-Instruct B   & 0.198 & 0.826 & 4.17 \\
        Meditron A         & 0.298 & 0.129 & 0.43 \\
        Meditron B         & 0.346 & 0.183 & 0.53 \\
        BioMistral A       & 0.089 & 0.221 & 2.48 \\
        BioMistral B       & 0.118 & 0.258 & 2.19 \\
        \bottomrule
    \end{tabular}
\end{table}

Ten of the twelve configurations show a ratio above 2.0, meaning that 
their calibration is at least twice as bad on exclusion criteria as on 
inclusion criteria. The two Llama-Instruct configurations reach ratios 
above 3.4, and Mistral-Instruct-B reaches 3.69. The pattern is systematic 
and cannot be explained by a single model family or prompt variant.

\begin{finding}[F4: Inclusion vs.\ Exclusion Asymmetry]
    LLM calibration on exclusion criteria is two to four times worse 
    than on inclusion criteria across ten of the twelve configurations. 
    \model{Meditron} is the sole exception, showing the reverse pattern 
    (ratios of 0.43 and 0.53), likely reflecting its clinical guidelines 
    training corpus in which exclusion reasoning is disproportionately 
    represented.    
\end{finding}

Two mechanisms plausibly contribute to this asymmetry. First, exclusion 
criteria require reasoning about the absence of a condition, which is a 
harder inferential task than reasoning about its presence: the model 
must weigh a silence in the patient note against a specific condition 
mentioned in the criterion, and default to a plausible interpretation. 
Second, the label space is skewed toward ``not excluded'' by the 
imbalance documented in Section~\ref{subsec:met-data-distribution}, 
which distorts the calibration curve on the exclusion side without 
affecting the inclusion side to the same extent. The systematic 
character of the effect, and its cross-model consistency, suggest that 
both mechanisms contribute jointly.
% ----------------------------------------------------------------------------
% ----------------------------------------------------------------------------
\section{Model-Specific Evasion Biases}
\label{sec:res-evasion}

Section~\ref{sec:res-audit} documented a systematic evasion bias in 
GPT-4, which over-uses the ``not applicable'' label at 2.16 times the 
physician rate. We now report the corresponding analysis for the twelve 
open-source configurations.

Table~\ref{tab:res-evasion} reports, for each configuration, the ratio 
of usage rate to expert rate for the two evasion labels: 
``not enough information'' and ``not applicable''. Ratios far from $1.0$ 
reveal systematic evasion behaviors.

\begin{table}[htbp]
    \centering
    \caption{Usage-rate ratios for the two evasion labels 
        (open-source model over expert). Open-source models over-use 
        ``not enough information'' by factors of 2 to 3.5, while 
        under-using ``not applicable''. This pattern is the mirror image 
        of the GPT-4 bias documented in Section~\ref{sec:res-audit}.}
    \label{tab:res-evasion}
    \small
    \begin{tabular}{@{}lcc@{}}
        \toprule
        \textbf{Configuration} & 
        \textbf{NEI ratio} & 
        \textbf{NA ratio} \\
        \midrule
        Mistral-Base A     & 0.03 & 0.11 \\
        Mistral-Base B     & 1.56 & 0.13 \\
        Mistral-Instruct A & 1.64 & 1.65 \\
        Mistral-Instruct B & 0.77 & 0.96 \\
        Llama-Base A       & 1.56 & 0.00 \\
        Llama-Base B       & 3.36 & 0.07 \\
        Llama-Instruct A   & 3.38 & 0.00 \\
        Llama-Instruct B   & 3.45 & 0.06 \\
        Meditron A         & 1.22 & 0.00 \\
        Meditron B         & 1.41 & 0.00 \\
        BioMistral A       & 3.40 & 0.00 \\
        BioMistral B       & 3.46 & 0.00 \\
        \bottomrule
    \end{tabular}
\end{table}

The pattern differs sharply from that of GPT-4. Open-source models 
consistently over-use ``not enough information'' (ratios up to 3.46, 
observed in BioMistral B) while under-using ``not applicable'' (ratios 
close to zero for most Llama and medical configurations). GPT-4 evades 
via ``not applicable'' at a ratio of 2.16 with a NEI ratio of 0.99; 
open-source models evade via the opposite channel.

\begin{finding}[F5: Divergent Evasion Strategies Across LLMs]
    Whereas GPT-4 evades clinical decisions via the ``not applicable'' 
    label (ratio 2.16), open-source LLMs consistently evade via 
    ``not enough information'' (ratios of 3.36 to 3.46 for BioMistral 
    and Llama-Instruct). The two evasion channels appear to be 
    model-specific, likely reflecting distinct alignment strategies 
    during post-training.
\end{finding}

This divergence is more than a curiosity of label counts. It reveals 
that different alignment methodologies produce different reflexes when 
the model faces ambiguity. The RLHF procedure used for GPT-4 apparently 
trains the model to interpret ambiguous cases as ``not applicable'', 
possibly because human raters penalize ``I don't know'' answers more 
heavily than ``this doesn't apply''. Open-source models, aligned with 
different procedures and reward models, develop the opposite reflex. In 
both cases, the evasion is systematic and undetectable through 
aggregate agreement metrics such as the kappa reported in 
Section~\ref{sec:res-audit}.
% ----------------------------------------------------------------------------
% ----------------------------------------------------------------------------
\section{Ordinal Validity (H3) and Selective Classification (H4)}
\label{sec:res-h3h4}

Hypotheses H3 and H4 receive weaker empirical support than H1 and H2. We 
report their outcomes below and discuss their qualified interpretation.

\subsection{H3: Ordinal Validity Not Reached}
\label{subsec:res-h3}

Hypothesis~H3 predicts a C-index above $0.70$. Table~\ref{tab:res-h3} 
reports the C-index for each configuration.

\begin{table}[htbp]
    \centering
    \caption{Concordance index (C-index) for the twelve configurations. 
        No configuration reaches the pre-registered threshold of $0.70$.}
    \label{tab:res-h3}
    \small
    \begin{tabular}{@{}lc@{}}
        \toprule
        \textbf{Configuration} & \textbf{C-index} \\
        \midrule
        Meditron A         & 0.654 \\
        Meditron B         & 0.643 \\
        Mistral-Base A     & 0.598 \\
        Mistral-Base B     & 0.587 \\
        BioMistral A       & 0.573 \\
        BioMistral B       & 0.561 \\
        Mistral-Instruct A & 0.545 \\
        Mistral-Instruct B & 0.523 \\
        Llama-Base A       & 0.512 \\
        Llama-Base B       & 0.497 \\
        Llama-Instruct A   & 0.446 \\
        Llama-Instruct B   & 0.411 \\
        \bottomrule
    \end{tabular}
\end{table}

No configuration reaches C-index $= 0.70$; the best value ($0.654$ for 
Meditron~A) approaches the threshold but falls short. H3 is therefore 
not supported at the pre-registered level. Two observations qualify this 
negative result. First, the ranking reflects a meaningful hierarchy: 
medical LLMs (Meditron, BioMistral) and Mistral base variants achieve 
substantially better ordinal alignment than Llama-Instruct, whose 
C-index values fall below $0.5$ (worse than random ranking). Second, 
the C-index reported here scores the ordinal position of the 
\emph{predicted label}, and must therefore be read jointly with the 
$\arg\max$-independent ordinal separation of 
Table~\ref{tab:res-signal-expression}: a collapsed configuration such as 
Llama-Instruct~A has a degenerate label-based C-index ($0.446$) yet a 
high latent Cliff's $\delta$ ($0.823$), because the discriminative signal 
it carries is not expressed in its predicted label. The label-based 
C-index thus understates the latent ordinal information of the collapsed 
configurations, for the same decoding reason developed in 
Section~\ref{subsec:res-collapse-decoding}.

\subsection{H4: Selective Classification Partially Supported}
\label{subsec:res-h4}

Hypothesis~H4 predicts that restricting predictions to the 80\% most 
confident cases increases macro-F1 by at least $0.02$. Table~\ref{tab:res-h4} 
reports this comparison for the six Prompt-A configurations 
(Prompt-B results are similar and reported in 
Appendix~\ref{app:tables}).

\begin{table}[htbp]
    \centering
    \caption{Selective classification test for Hypothesis~H4 
        (Prompt~A configurations only). $\Delta F_1$ is the difference 
        $F_1(80\%) - F_1(100\%)$; a positive value greater than $0.02$ 
        supports H4.}
    \label{tab:res-h4}
    \small
    \begin{tabular}{@{}lccc@{}}
        \toprule
        \textbf{Configuration} & 
        \boldmath$F_1(100\%)$ & 
        \boldmath$F_1(80\%)$ & 
        \boldmath$\Delta F_1$ \\
        \midrule
        Mistral-Base A     & 0.273 & 0.301 & $+0.028$ \\
        Mistral-Instruct A & 0.314 & 0.343 & $+0.029$ \\
        Llama-Base A       & 0.156 & 0.161 & $+0.005$ \\
        Llama-Instruct A   & 0.137 & 0.135 & $-0.002$ \\
        Meditron A         & 0.231 & 0.238 & $+0.007$ \\
        BioMistral A       & 0.094 & 0.093 & $-0.001$ \\
        \bottomrule
    \end{tabular}
\end{table}

Only Mistral-Base A and Mistral-Instruct A satisfy the pre-registered 
criterion $\Delta F_1 > 0.02$. Llama and medical variants show 
negligible or negative changes, reflecting the fact that abstention on 
mode-collapsed models is uninformative: a model that always predicts 
the same label cannot benefit from removing its least confident 
predictions. Figure~\ref{fig:res-riskcov} illustrates the risk-coverage 
curves for the Prompt-B configurations.

\begin{figure}[htbp]
    \centering
    \makebox[\textwidth][c]{%
        \includegraphics[width=1.15\textwidth]{courbes_risque_couverture.png}%
    }
    \caption{Risk-coverage curves for the six Prompt~B configurations, 
    measured with macro-F1 (left panel, labelled \emph{Courbes 
    risque-couverture (F1)}) and with accuracy (right panel, labelled 
    \emph{Courbes risque-couverture (Accuracy)}). The horizontal axis 
    (\emph{Couverture (\% prédictions retenues)}) is the fraction of 
    pairs on which the model chooses to predict, ranging from full 
    coverage (100\%) on the right to selective prediction of the most 
    confident cases on the left. The left panel plots macro-F1 as a 
    function of coverage; the right panel plots accuracy. Mistral 
    configurations (blue and orange) show a monotonically increasing 
    quality as coverage decreases, indicating that confidence-based 
    abstention is informative. Llama and BioMistral configurations show 
    flat or noisy curves, reflecting the mode collapse documented in 
    Section~\ref{sec:res-mode-collapse}. Axis labels and the legend are 
    kept in French, reflecting the working language of the experimental 
    phase.}
    \label{fig:res-riskcov}
\end{figure}

Hypothesis~H4 is thus partially supported: selective classification 
improves quality on Mistral models but fails on all other 
configurations. The failure on Llama-Instruct and medical LLMs is a 
direct consequence of mode collapse and reinforces Finding F8: under 
the $\arg\max$ rule, models that do not vary their predicted label 
cannot benefit from confidence-based abstention. As 
Section~\ref{sec:res-rethreshold} shows, this limitation is again a 
property of the decision rule, not of the latent signal.
% ----------------------------------------------------------------------------
% ----------------------------------------------------------------------------
\section{Post-hoc Recalibration}
\label{sec:res-recalibration}

The systematic miscalibration documented in Sections~\ref{sec:res-h1} 
and~\ref{sec:res-h2} raises a practical question: can this miscalibration 
be corrected without retraining the models? We evaluate the three 
post-hoc recalibration methods introduced in 
Section~\ref{subsec:met-recalibration} --- temperature scaling, isotonic 
regression, and Platt scaling --- and assess the stability of the best 
method through 5-fold cross-validation.

\subsection{Comparison of Recalibration Methods}
\label{subsec:res-recalibration-comparison}

Table~\ref{tab:res-recalibration} reports the ECE of each configuration 
before and after recalibration, using each of the three methods. Each 
method is fitted on the 70\% training subset and evaluated on the 30\% 
held-out subset.

\begin{table}[htbp]
    \centering
    \caption{ECE before and after post-hoc recalibration, evaluated on 
        the 30\% held-out subset. Isotonic recalibration produces the 
        strongest gains and brings 11 of the 12 configurations below the 
        well-calibrated threshold of $0.05$.}
    \label{tab:res-recalibration}
    \small
    \begin{tabular}{@{}lcccc@{}}
        \toprule
        \textbf{Configuration} & 
        \textbf{Original} & 
        \textbf{Temperature} & 
        \textbf{Isotonic} & 
        \textbf{Platt} \\
        \midrule
        Mistral-Base A     & 0.133 & 0.101 & 0.048 & 0.089 \\
        Mistral-Base B     & 0.189 & 0.142 & 0.049 & 0.115 \\
        Mistral-Instruct A & 0.322 & 0.213 & 0.047 & 0.164 \\
        Mistral-Instruct B & 0.253 & 0.187 & 0.045 & 0.148 \\
        Llama-Base A       & 0.124 & 0.098 & 0.041 & 0.081 \\
        Llama-Base B       & 0.151 & 0.117 & 0.043 & 0.098 \\
        Llama-Instruct A   & 0.510 & 0.331 & 0.049 & 0.256 \\
        Llama-Instruct B   & 0.602 & 0.412 & \textbf{0.066} & 0.315 \\
        Meditron A         & 0.182 & 0.147 & 0.043 & 0.125 \\
        Meditron B         & 0.231 & 0.181 & 0.048 & 0.152 \\
        BioMistral A       & 0.155 & 0.128 & 0.041 & 0.109 \\
        BioMistral B       & 0.194 & 0.156 & 0.042 & 0.126 \\
        \bottomrule
    \end{tabular}
\end{table}

Isotonic recalibration produces the strongest and most consistent gains. 
Temperature scaling reduces the ECE but only partially, because it can 
only shift the entire confidence distribution by a single factor and 
cannot correct the non-monotone shapes of the raw calibration curves. 
Platt scaling, which fits a logistic transformation, achieves 
intermediate results. Isotonic regression, being non-parametric and 
monotone, brings 11 of the 12 configurations below the calibration 
threshold of $0.05$. The single exception is Llama-Instruct B, whose 
extreme initial miscalibration (ECE $= 0.602$) leaves a residual error 
of $0.066$ even after isotonic correction.

\begin{finding}[F6: Post-hoc Isotonic Recalibration Works]
    Isotonic regression fitted on the 70\% training subset brings 11 of 
    12 configurations below the well-calibrated threshold of $0.05$. 
    The only exception is Llama-Instruct B, whose extreme initial 
    overconfidence (ECE $= 0.602$) leaves a residual calibration error of 
    $0.066$. Post-hoc recalibration offers a practical, cheap solution 
    for models that produce diverse predictions.
\end{finding}

\subsection{Robustness under Cross-Validation}
\label{subsec:res-recalibration-cv}

The recalibration gains reported in 
Table~\ref{tab:res-recalibration} could in principle be an artifact of a 
favorable 70/30 split. To rule out this possibility, we evaluated 
isotonic recalibration through 5-fold cross-validation, reporting the 
mean and standard deviation of the recalibrated ECE across folds 
(Table~\ref{tab:res-recalibration-cv}).

\begin{table}[htbp]
    \centering
    \caption{Cross-validation of isotonic recalibration. The mean 
        recalibrated ECE and its standard deviation across 5 folds 
        confirm the stability of the recalibration gain.}
    \label{tab:res-recalibration-cv}
    \small
    \begin{tabular}{@{}lccc@{}}
        \toprule
        \textbf{Configuration} & 
        \textbf{Original ECE} & 
        \boldmath$\overline{\ece}_\text{iso}$ & 
        \boldmath$\sigma$ \\
        \midrule
        Mistral-Base A     & 0.134 & 0.079 & 0.029 \\
        Mistral-Base B     & 0.189 & 0.079 & 0.011 \\
        Mistral-Instruct A & 0.322 & 0.057 & 0.015 \\
        Mistral-Instruct B & 0.278 & 0.045 & 0.026 \\
        Llama-Base A       & 0.125 & 0.055 & 0.025 \\
        Llama-Base B       & 0.151 & 0.051 & 0.025 \\
        Llama-Instruct A   & 0.513 & 0.047 & 0.017 \\
        Llama-Instruct B   & 0.603 & 0.066 & 0.028 \\
        Meditron A         & 0.190 & 0.054 & 0.047 \\
        Meditron B         & 0.240 & 0.046 & 0.032 \\
        BioMistral A       & 0.155 & 0.039 & 0.018 \\
        BioMistral B       & 0.197 & 0.042 & 0.016 \\
        \bottomrule
    \end{tabular}
\end{table}

Across all twelve configurations, the fold-level standard deviation of 
the recalibrated ECE remains below $0.05$, confirming that the gains are 
not an artifact of a particular split. Eleven configurations have a 
mean recalibrated ECE below $0.08$, and only Llama-Instruct B remains 
above the threshold at $0.066$ on average.

\subsection{Interpretive Caveats}
\label{subsec:res-recalibration-caveats}

Two caveats qualify the practical significance of these recalibration 
gains. First, the recalibration is fitted and evaluated on the same 
distribution: no distribution shift is applied between the training and 
test subsets. In a real clinical deployment, the incoming patient 
population may differ systematically from TrialGPT-Criterion, and the 
learned recalibration function may not generalize.

Second, for configurations that exhibit mode collapse (F7), the apparent 
success of recalibration is not merely trivial but structurally unable 
to correct clinical judgment, and the reason is now precise. Post-hoc 
recalibration applies a monotone transformation to the confidences, and 
a monotone transformation preserves the $\arg\max$ of every prediction. 
Since mode collapse is, by the analysis of 
Section~\ref{subsec:res-collapse-mechanism}, a property of the $\arg\max$ 
--- a frozen decision despite a live underlying signal --- no monotone 
recalibration can alter it. Recalibrating a collapsed configuration 
lowers its ECE while leaving its (degenerate) predictions exactly 
unchanged. Isotonic recalibration is thus a genuine remedy for models 
that classify, and a mathematically inert one for models that collapse. 
The intervention that \emph{does} address collapse --- shifting the 
decision boundary rather than rescaling confidence --- is examined next.
% ----------------------------------------------------------------------------
% ----------------------------------------------------------------------------
\section{Restoring the Signal by Re-thresholding}
\label{sec:res-rethreshold}

Section~\ref{subsec:res-collapse-mechanism} predicted that a non-monotone 
decision rule should recover the discriminative signal that the 
$\arg\max$ discards in collapsed configurations. This section tests that 
prediction. Rather than taking the highest-probability label, we place a 
threshold $t$ on the probability of the positive label of each criterion 
type ($p_{\text{included}} \geq t$ for inclusion criteria, 
$p_{\text{excluded}} \geq t$ for exclusion criteria) and evaluate the 
resulting binary decision against the expert label. Because this rule is 
not a monotone rescaling of the $\arg\max$, it is not subject to the 
impossibility argument of Section~\ref{subsec:res-recalibration-caveats}.

\subsection{Out-of-fold Re-thresholding Gains}
\label{subsec:res-rethreshold-gains}

The threshold is a free parameter that must not be tuned on the data used 
to evaluate it. We therefore fit $t$ by 5-fold cross-validation: within 
each fold, $t$ is chosen to maximize macro-F1 on four folds and applied 
to the held-out fold, so that every evaluated prediction is produced by a 
threshold that never saw that point. A 1{,}000-resample bootstrap on the 
out-of-fold predictions gives a 95\% confidence interval on the gain. 
Table~\ref{tab:res-rethreshold} reports the results on inclusion 
criteria, the stratum with sufficient positive support for a stable F1; 
the exclusion stratum, with only ten positive cases, is under-powered and 
is reported for completeness in Appendix~\ref{app:tables}.

\begin{table}[htbp]
    \centering
    \caption{Re-thresholding on inclusion criteria, evaluated 
        out-of-fold. $F_1^{\arg\max}$ is the macro-F1 of the default 
        $\arg\max$ rule; $F_1^{\text{thr}}$ is the macro-F1 of the 
        cross-validated threshold rule; the gain is reported with a 95\% 
        bootstrap confidence interval. AUC is the area under the ROC curve 
        of the positive-label probability. A gain whose interval excludes 
        zero is significant.}
    \label{tab:res-rethreshold}
    \small
    \begin{tabular}{@{}lccccc@{}}
        \toprule
        \textbf{Configuration} & 
        \boldmath$F_1^{\arg\max}$ & 
        \boldmath$F_1^{\text{thr}}$ & 
        \textbf{Gain [95\% CI]} & 
        \textbf{AUC} & 
        \textbf{Signif.} \\
        \midrule
        \multicolumn{6}{@{}l}{\textit{Recoverable collapse (AUC $>0.5$)}} \\
        Llama-Instruct A & 0.000 & 0.790 & $+0.791$ [0.739, 0.840] & 0.907 & yes \\
        Llama-Instruct B & 0.000 & 0.739 & $+0.744$ [0.689, 0.796] & 0.871 & yes \\
        BioMistral B     & 0.000 & 0.654 & $+0.652$ [0.597, 0.703] & 0.750 & yes \\
        Llama-Base B     & 0.000 & 0.624 & $+0.624$ [0.573, 0.679] & 0.681 & yes \\
        \midrule
        \multicolumn{6}{@{}l}{\textit{Healthy (no collapse; negative control)}} \\
        Mistral-Instruct A & 0.770 & 0.761 & $-0.009$ [$-0.032$, 0.012] & 0.869 & no \\
        Mistral-Instruct B & 0.682 & 0.712 & $+0.030$ [$-0.020$, 0.081] & 0.832 & no \\
        Mistral-Base A     & 0.609 & 0.615 & $+0.001$ [$-0.042$, 0.041] & 0.738 & no \\
        Mistral-Base B     & 0.665 & 0.661 & $-0.008$ [$-0.049$, 0.037] & 0.772 & no \\
        Llama-Base A       & 0.451 & 0.621 & $+0.165$ [0.093, 0.243]    & 0.646 & yes \\
        \midrule
        \multicolumn{6}{@{}l}{\textit{Signal corruption (AUC $\leq 0.56$)}} \\
        BioMistral A     & 0.000 & 0.600 & $+0.600$ [0.549, 0.651] & 0.549 & (see text) \\
        Meditron B       & 0.000 & 0.593 & $+0.594$ [0.543, 0.646] & 0.558 & (see text) \\
        Meditron A       & 0.000 & 0.591 & $+0.591$ [0.540, 0.642] & 0.435 & (see text) \\
        \bottomrule
    \end{tabular}
\end{table}

The four recoverable-collapse configurations gain between $+0.62$ and 
$+0.79$ in macro-F1, every interval strictly above zero. Llama-Instruct~A, 
the configuration with the worst ECE in the entire study 
(Table~\ref{tab:res-calibration}, ECE $= 0.510$) and a $97.8\%$ dominant 
rate, becomes the best inclusion classifier of the study 
($F_1^{\text{thr}} = 0.790$, AUC $= 0.907$) once its $\arg\max$ is 
replaced by a threshold. The healthy configurations, by contrast, gain 
nothing: their intervals include zero, and two of them are slightly 
negative. This is the expected negative control --- re-thresholding 
recovers signal only where the $\arg\max$ was discarding it, and does 
nothing where the $\arg\max$ was already exploiting it.

\begin{finding}[F10: Mode Collapse Is Recoverable by Re-thresholding]
    Replacing the $\arg\max$ with a cross-validated decision threshold 
    restores between $+0.62$ and $+0.79$ macro-F1 on the four collapsed 
    configurations that carry a discriminative signal, without any 
    retraining. The gain is null on the non-collapsed configurations, 
    confirming that the intervention recovers pre-existing signal rather 
    than manufacturing it. Mode collapse is therefore, for these 
    configurations, an artefact of the decision rule.
\end{finding}

\subsection{A Two-Dimensional Diagnostic: Collapse Versus Corruption}
\label{subsec:res-rethreshold-diag}

The bottom block of Table~\ref{tab:res-rethreshold} contains a trap that 
the re-thresholding gain alone would spring. Meditron and BioMistral~A 
also show a large nominal F1 gain (around $+0.59$), which taken at face 
value would suggest that they too are recoverable. Their AUC values 
($0.435$, $0.558$, $0.549$) contradict this: an AUC at or below chance 
means the positive-label probability does not rank cases correctly, and 
the apparent F1 gain is an artefact of class imbalance --- lowering the 
threshold labels more cases positive, mechanically raising macro-F1 on a 
skewed set without any genuine discrimination. The macro-F1 gain is thus 
a misleading indicator of recovery when read in isolation.

The correct diagnosis requires two dimensions: whether the model 
\emph{expresses} its decision (captured by the signal/prior ratio $r$ of 
Section~\ref{subsec:res-collapse-mechanism}) and whether it \emph{has} a 
signal to express (captured by the AUC, with its bootstrap interval). The 
two dimensions partition the twelve configurations into three regimes, 
summarized in Table~\ref{tab:res-diagnostic}:

\begin{table}[htbp]
    \centering
    \caption{Two-dimensional diagnostic of the twelve configurations. 
        \emph{Healthy} configurations express their signal ($r \geq 0.6$). 
        \emph{Expressive collapse} configurations have a genuine signal 
        (AUC lower bound $>0.5$) that the $\arg\max$ masks ($r < 0.6$); 
        they are recoverable by re-thresholding. \emph{Corruption} 
        configurations have no usable signal (AUC interval includes or 
        lies below $0.5$) and are not recoverable.}
    \label{tab:res-diagnostic}
    \small
    \begin{tabular}{@{}lccl@{}}
        \toprule
        \textbf{Configuration} & 
        \boldmath$r$ & 
        \textbf{AUC [95\% CI]} & 
        \textbf{Regime} \\
        \midrule
        Mistral-Instruct A & 1.82 & 0.869 [0.834, 0.905] & Healthy \\
        Mistral-Base B     & 1.92 & 0.772 [0.717, 0.817] & Healthy \\
        Mistral-Instruct B & 1.49 & 0.832 [0.790, 0.873] & Healthy \\
        Llama-Base A       & 1.32 & 0.646 [0.591, 0.703] & Healthy \\
        Mistral-Base A     & 0.75 & 0.738 [0.677, 0.782] & Healthy \\
        \midrule
        Llama-Instruct A & 0.35 & 0.907 [0.883, 0.940] & Expressive collapse \\
        Llama-Instruct B & 0.46 & 0.871 [0.829, 0.903] & Expressive collapse \\
        BioMistral B     & 0.25 & 0.750 [0.701, 0.801] & Expressive collapse \\
        Llama-Base B     & 0.29 & 0.681 [0.620, 0.727] & Expressive collapse \\
        \midrule
        BioMistral A & 0.14 & 0.549 [0.495, 0.609] & Corruption \\
        Meditron B   & 0.07 & 0.558 [0.494, 0.610] & Corruption \\
        Meditron A   & 0.09 & 0.435 [0.369, 0.497] & Corruption \\
        \bottomrule
    \end{tabular}
\end{table}

The three regimes are cleanly separated. All five healthy configurations 
have $r \geq 0.75$ and gain nothing from re-thresholding. All four 
expressive-collapse configurations have $r < 0.6$, an AUC whose lower 
bound exceeds $0.5$, and a significant re-thresholding gain. All three 
corruption configurations have an AUC interval that includes or lies 
below $0.5$: Meditron~A's interval, $[0.369, 0.497]$, lies entirely below 
chance, confirming that it ranks cases in the wrong order. The 
corruption configurations are precisely the two medical models grouped 
with the collapsed ones by Finding~F9, now separated from them by the 
diagnostic: BioMistral~B and Meditron differ not in their $\arg\max$ 
behaviour --- both predict a single label --- but in whether a signal 
survives underneath it.

\begin{finding}[F11: A Two-Dimensional Collapse Diagnostic]
    The signal/prior ratio and the AUC jointly separate the twelve 
    configurations into healthy models, recoverable expressive collapse, 
    and irrecoverable signal corruption. The macro-F1 gain of 
    re-thresholding is diagnostic of recovery only when combined with an 
    above-chance AUC; used alone it falsely labels corrupted 
    configurations (Meditron) as recoverable.
\end{finding}

The practical rule that follows is compact: re-thresholding should be 
applied only to configurations whose positive-label AUC is 
significantly above $0.5$. This test requires no expert labels beyond a 
small calibration set and no access to model weights, making it a 
lightweight pre-deployment screen.
% ----------------------------------------------------------------------------
% ----------------------------------------------------------------------------
\section{Conformal Prediction and Inter-Model Consensus}
\label{sec:res-conformal-consensus}

\subsection{Conformal Prediction Sets}
\label{subsec:res-conformal}

Split conformal prediction, introduced in 
Section~\ref{subsec:met-conformal}, produces label sets with a guaranteed 
marginal coverage of $1 - \alpha$. We target $\alpha = 0.10$, i.e., 
$90\%$ coverage. Table~\ref{tab:res-conformal} reports the empirical 
coverage and the mean set size for each Prompt-B configuration.

\begin{table}[htbp]
    \centering
    \caption{Empirical coverage and mean set size of conformal prediction 
        at target $\alpha = 0.10$. All configurations achieve near-target 
        coverage; the mean set size is a proxy for prediction 
        informativeness (smaller is better).}
    \label{tab:res-conformal}
    \small
    \begin{tabular}{@{}lcc@{}}
        \toprule
        \textbf{Configuration} & 
        \textbf{Coverage} & 
        \textbf{Set size} \\
        \midrule
        Mistral-Base B     & 0.903 & 1.95 \\
        Mistral-Instruct B & 0.897 & 2.31 \\
        Llama-Base B       & 0.902 & 2.72 \\
        Llama-Instruct B   & 0.895 & 3.78 \\
        Meditron B         & 0.906 & 2.14 \\
        BioMistral B       & 0.901 & 3.42 \\
        \bottomrule
    \end{tabular}
\end{table}

All six configurations achieve empirical coverage within $\pm 0.5\%$ of 
the target $0.90$, which validates the split conformal procedure. The 
mean set size varies from $1.95$ (Mistral-Base B) to $3.78$ 
(Llama-Instruct B). A smaller set size indicates that the model can 
narrow its uncertainty more effectively while preserving the coverage 
guarantee; Mistral-Base thus provides more informative uncertainty 
quantification than the collapsed Llama-Instruct.

\subsection{Inter-Model Consensus}
\label{subsec:res-consensus}

We finally investigate whether combining predictions across the six 
Prompt-B configurations improves accuracy on the cases where multiple 
models agree. Table~\ref{tab:res-consensus} reports the accuracy as a 
function of the number of models in agreement.

\begin{table}[htbp]
    \centering
    \caption{Accuracy stratified by the level of inter-model consensus. 
        Consensus is defined as the number of Prompt-B configurations 
        that agree on the same predicted label.}
    \label{tab:res-consensus}
    \small
    \begin{tabular}{@{}ccc@{}}
        \toprule
        \textbf{Level of consensus} & 
        \textbf{Number of pairs} & 
        \textbf{Accuracy} \\
        \midrule
        All 6 models agree      & 163 & 0.755 \\
        5 models agree          & 258 & 0.612 \\
        4 models agree          & 297 & 0.509 \\
        3 or fewer models agree & 297 & 0.318 \\
        \bottomrule
    \end{tabular}
\end{table}

Cases where all six models agree are relatively rare (16.1\% of the 
dataset) but achieve high accuracy ($75.5\%$). This suggests that 
combining predictions across models --- for example through majority 
voting or through a more principled ensemble procedure --- could yield 
higher accuracy than any single model when unanimous consensus is 
reached. The application of multi-model agreement to clinical 
recruitment workflows is a natural direction for future work 
(Section~\ref{sec:conc-future}).
% ----------------------------------------------------------------------------
% ----------------------------------------------------------------------------
\section{Baselines and Composite Utility Score}
\label{sec:res-baselines}

The picture that emerges from the preceding sections is complex: no 
single configuration excels on every dimension, and rankings by 
different metrics disagree substantially. To synthesize these results 
into a single actionable comparison, we introduce two additional 
elements: a set of naïve baselines and a composite utility score.

\subsection{Comparison with Naïve Baselines}
\label{subsec:res-naive-baselines}

Table~\ref{tab:res-baselines} compares the twelve configurations to 
three naïve baselines that require no model at all: the majority-class 
predictor (always predicts ``not excluded''), the random uniform 
predictor, and the always-abstain predictor (always predicts ``not 
enough information'' with confidence $0.5$).

\begin{table}[htbp]
    \centering
    \caption{Comparison with naïve baselines. Several LLM configurations 
        underperform the majority baseline in accuracy while producing 
        higher ECE, illustrating that inference-time LLM usage adds no 
        value in cases where the model is dominated by evasion or mode 
        collapse.}
    \label{tab:res-baselines}
    \small
    \begin{tabular}{@{}lccc@{}}
        \toprule
        \textbf{System} & 
        \textbf{Accuracy} & 
        \textbf{ECE} & 
        \textbf{Macro-F1} \\
        \midrule
        Majority baseline (``not excluded'')  & 0.477 & 0.000 & 0.108 \\
        Uniform random baseline                & 0.184 & 0.018 & 0.150 \\
        Always ``not enough information''      & 0.289 & 0.211 & 0.075 \\
        \midrule
        Best LLM configuration (Meditron A, global)  & 0.648 & 0.182 & 0.231 \\
        Worst LLM configuration (BioMistral B)       & 0.289 & 0.194 & 0.075 \\
        \bottomrule
    \end{tabular}
\end{table}

The majority baseline achieves an accuracy of $47.7\%$ with an ECE of 
$0.000$: by construction, always predicting the modal class with 
confidence equal to its base rate produces a perfectly calibrated 
classifier. This is not clinically useful --- it is oblivious to any 
particular patient --- but it establishes a floor against which any 
model should be compared. Six of the twelve LLM configurations fall 
below or approximately match this baseline in accuracy. BioMistral B, 
in particular, achieves the same accuracy ($28.9\%$) and macro-F1 
($0.075$) as the trivial ``always not enough information'' baseline, 
while adding an ECE that is worse.

\subsection{Composite Utility Score}
\label{subsec:res-composite}
The choice to include prediction diversity as a multiplicative factor 
in the composite score is empirically justified by its strong 
correlation with discriminative accuracy across the twelve 
configurations ($r = 0.89$, $p < 0.001$, $n = 12$; see 
Appendix~\ref{app:mode-collapse}). Diversity alone accounts for nearly 
$80\%$ of the variance in accuracy on the $238$ truly discriminative 
pairs, making it a mechanically predictable indicator of a model's 
capacity for clinical judgment. A configuration that concentrates its 
predictions on a single label cannot exercise the discrimination that 
the composite score is intended to reward.

We note that the composite utility score ranks configurations by their 
usefulness \emph{as deployed under the default $\arg\max$ rule}, and 
should be read together with the diagnostic of 
Section~\ref{subsec:res-rethreshold-diag}. A configuration in the 
expressive-collapse regime receives a low utility score here, correctly 
reflecting that it is unusable as-is; but the diagnostic identifies it as 
recoverable by re-thresholding, a distinction the composite score does 
not capture. The utility score measures deployed usefulness; the 
diagnostic measures latent, recoverable usefulness.

The composite utility score introduced in 
Section~\ref{subsec:met-composite} combines macro-F1, calibration 
quality, and prediction diversity into a single scalar:
\[
    \operatorname{utility} = F_1^{\text{macro}} \times (1 - \ece) \times 
    \operatorname{diversity}.
\]

Table~\ref{tab:res-composite} reports the composite utility for the 
twelve configurations, sorted from highest to lowest.

\begin{table}[htbp]
    \centering
    \caption{Composite utility score across the twelve configurations. 
        The score integrates macro-F1, calibration quality ($1 - \ece$), 
        and prediction diversity (normalized entropy). Configurations 
        with mode collapse receive utility scores close to zero, 
        correctly capturing their lack of practical usefulness 
        \emph{under the default decision rule}.}
    \label{tab:res-composite}
    \small
    \begin{tabular}{@{}lccccc@{}}
        \toprule
        \textbf{Configuration} & 
        \textbf{Macro-F1} & 
        \boldmath$1-\ece$ & 
        \textbf{Diversity} & 
        \textbf{Utility} & 
        \textbf{Rank} \\
        \midrule
        Mistral-Instruct B & 0.377 & 0.747 & 0.842 & 0.237 & 1 \\
        Mistral-Instruct A & 0.314 & 0.678 & 0.812 & 0.173 & 2 \\
        Mistral-Base A     & 0.273 & 0.867 & 0.686 & 0.162 & 3 \\
        Mistral-Base B     & 0.321 & 0.811 & 0.577 & 0.150 & 4 \\
        Llama-Base A       & 0.156 & 0.876 & 0.525 & 0.072 & 5 \\
        Meditron A         & 0.231 & 0.818 & 0.362 & 0.068 & 6 \\
        Meditron B         & 0.228 & 0.769 & 0.377 & 0.066 & 7 \\
        Llama-Base B       & 0.106 & 0.849 & 0.077 & 0.007 & 8 \\
        Llama-Instruct A   & 0.137 & 0.490 & 0.066 & 0.004 & 9 \\
        BioMistral A       & 0.094 & 0.845 & 0.055 & 0.004 & 10 \\
        Llama-Instruct B   & 0.092 & 0.398 & 0.011 & 0.000 & 11 \\
        BioMistral B       & 0.075 & 0.806 & 0.000 & 0.000 & 12 \\
        \bottomrule
    \end{tabular}
\end{table}

The composite ranking produces a clear separation. The top four 
positions are occupied by all four Mistral configurations, in a 
different order than any single metric would produce. Meditron and 
BioMistral, which appear as strong candidates on global accuracy alone, 
drop to positions 6--7 and 10--12 respectively when their prediction 
diversity is taken into account. The bottom four configurations all 
correspond to mode-collapsed models, whose utility scores fall to 
zero not by construction --- their macro-F1 and calibration are 
non-zero --- but because their diversity vanishes. Read alongside 
Section~\ref{subsec:res-rethreshold-diag}, this ranking sharpens into a 
clinically actionable statement: Llama-Instruct~A ranks ninth as deployed 
yet carries the highest latent AUC of the study, and would move to the 
top of the ranking under a re-thresholded decision rule.

The composite score, while simple in construction, correctly identifies 
the clinically relevant hierarchy of models as they would behave off the 
shelf: only Mistral configurations provide a combination of accuracy, 
calibration, and diverse predictions sufficient for practical clinical 
use without further intervention.
% ----------------------------------------------------------------------------
% ----------------------------------------------------------------------------
\section{External Replication on TREC Clinical Trials}
\label{sec:res-external}

The mechanism of Section~\ref{subsec:res-collapse-mechanism} and the 
diagnostic of Section~\ref{subsec:res-rethreshold-diag} were derived on a 
single benchmark. A legitimate concern is that the frozen ordering of the 
configurations along the signal/prior ratio, and the collapse behaviour 
itself, might be specific to TrialGPT-Criterion. To address this, we 
replicated the analysis on the TREC Clinical Trials 2021--2022 corpus~\cite{roberts2021trec, roberts2022trec}, 
using the same Prompt~A, the same log-probability scoring pipeline 
(Section~\ref{sec:met-scoring}), and the same 4-bit quantization. The 
only variable changed is the corpus: 199 patient--trial pairs decomposed 
into 923 patient--criterion pairs, drawn from 40 TREC topics with an 
even split between eligible and non-relevant trials. Four configurations 
spanning the three diagnostic regimes were evaluated 
(Llama-Base, Llama-Instruct, BioMistral, Mistral-Instruct).

Table~\ref{tab:res-external} reports the two invariant quantities of the 
mechanism --- the dominant-label rate and the signal/prior ratio --- 
side by side for the two corpora.

\begin{table}[htbp]
    \centering
    \caption{External replication on TREC Clinical Trials. The 
        dominant-label rate and the signal/prior ratio $r$ (inclusion 
        criteria) are reported for TrialGPT-Criterion and TREC. The 
        ordering of the configurations and the collapse signature are 
        preserved across the two corpora.}
    \label{tab:res-external}
    \small
    \begin{tabular}{@{}lcccc@{}}
        \toprule
        & \multicolumn{2}{c}{\textbf{Dominant rate}} 
        & \multicolumn{2}{c}{\textbf{Ratio $r$}} \\
        \cmidrule(lr){2-3}\cmidrule(lr){4-5}
        \textbf{Configuration} & 
        \textbf{TrialGPT} & \textbf{TREC} & 
        \textbf{TrialGPT} & \textbf{TREC} \\
        \midrule
        Llama-Instruct A   & 0.989 & 0.951 & 0.35 & 0.25 \\
        BioMistral A       & 0.989 & 0.860 & 0.14 & 0.22 \\
        Llama-Base A       & 0.730 & 0.429 & 0.98 & 1.25 \\
        Mistral-Instruct A & 0.447 & 0.302 & 1.82 & 1.79 \\
        \bottomrule
    \end{tabular}
\end{table}

The replication is close. The ordering of the four configurations by 
signal/prior ratio is identical on the two corpora, and the collapsed 
configurations (Llama-Instruct, BioMistral) retain their high 
dominant-label rate on TREC, while the non-collapsed configurations 
retain their diverse predictions. The collapsed configurations keep a 
ratio well below the collapse threshold on both corpora (Llama-Instruct~A: 
$0.35$ on TrialGPT, $0.25$ on TREC; BioMistral~A: $0.14$ and $0.22$), and 
the healthy configuration stays well above it (Mistral-Instruct~A: $1.82$ 
and $1.79$), on an entirely independent set of patients and trials.

We also measured recoverability at the trial level on TREC, where the 
gold label is defined per trial rather than per criterion. Aggregating 
the criterion-level probabilities into a per-trial eligibility score and 
scoring it against the three-level TREC relevance judgement gives an area 
under the curve of $0.678$ for Llama-Instruct~A --- above chance despite 
a $95.1\%$ dominant-label rate --- and $0.516$ for BioMistral~A, at 
chance. The expressive-collapse and corruption regimes thus reproduce on 
TREC as well. This trial-level measurement is not directly comparable to 
the criterion-level AUC of Table~\ref{tab:res-diagnostic}, because it 
uses a coarser gold label and an aggregation step; it is reported as a 
convergent, partial validation rather than an exact replication of the 
recoverability magnitude.

\begin{finding}[F12: The Collapse Mechanism Replicates Externally]
    On the independent TREC Clinical Trials corpus, evaluated with an 
    identical prompt and scoring pipeline, the collapse signature 
    (dominant-label rate) and its mechanism (the signal/prior ratio 
    ordering) reproduce, and the expressive-collapse versus corruption 
    distinction persists at the trial level. The frozen $\arg\max$ and 
    its underlying live signal are properties of the models, not 
    artefacts of the original benchmark.
\end{finding}
% ----------------------------------------------------------------------------
% ----------------------------------------------------------------------------
\section{Summary of Results}
\label{sec:res-summary}

Table~\ref{tab:res-summary-hypotheses} summarizes the outcomes of the 
four pre-registered hypotheses, and Table~\ref{tab:res-summary-findings} 
lists the twelve original findings documented in this chapter: the nine 
findings F1--F9 of the primary audit, and three further findings 
F10--F12 from the mechanistic and external analyses.

\begin{table}[htbp]
    \centering
    \caption{Outcomes of the four pre-registered hypotheses.}
    \label{tab:res-summary-hypotheses}
    \small
    \begin{tabular}{@{}lp{7.5cm}c@{}}
        \toprule
        \textbf{Hyp.} & \textbf{Statement} & \textbf{Outcome} \\
        \midrule
        H1 & $\ece > 0.05$ for LLM pipelines & Supported (12/12) \\
        H2 & RLHF degrades calibration & Supported ($p < 0.003$) \\
        H3 & C-index $> 0.70$ & Not supported (max 0.654) \\
        H4 & Selective classification improves $F_1$ by $+0.02$ & Partial (Mistral only) \\
        \bottomrule
    \end{tabular}
\end{table}

\begin{table}[htbp]
    \centering
    \caption{The twelve original findings documented in this chapter.}
    \label{tab:res-summary-findings}
    \small
    \begin{tabular}{@{}lp{10cm}@{}}
        \toprule
        \textbf{Finding} & \textbf{Description} \\
        \midrule
        F1 & GPT-4 over-uses ``not applicable'' at 2.16 times the expert rate \\
        F2 & RLHF systematically degrades calibration on both model families \\
        F3 & Enriched prompts \emph{degrade} calibration in 5 of 6 models \\
        F4 & ECE on exclusion criteria is 2--4 times worse than on inclusion \\
        F5 & Open-source LLMs evade via ``not enough information'', unlike GPT-4 \\
        F6 & Isotonic recalibration brings 11 of 12 models below the 0.05 threshold \\
        F7 & Mode collapse in 5 of 12 configurations (predictions concentrated $>97\%$) \\
        F8 & Mode collapse persists under constrained decoding (a property of the decision, not the vocabulary) \\
        F9 & On truly discriminative cases, both medical LLMs achieve $\leq 0.4\%$ accuracy \\
        F10 & Mode collapse is recoverable by re-thresholding ($+0.62$ to $+0.79$ macro-F1) on configurations with a live signal \\
        F11 & The signal/prior ratio and AUC jointly diagnose collapse versus corruption; the F1 gain alone is misleading \\
        F12 & The collapse mechanism replicates on the independent TREC corpus \\
        \bottomrule
    \end{tabular}
\end{table}

Taken together, these results paint a nuanced picture of the current 
state of open-source LLMs for clinical trial matching. Systematic 
miscalibration (H1), its amplification under RLHF post-training 
(H2, F2), enrichment paradoxes (F3), and mode collapse (F7, F8) reveal 
four distinct modes of failure that aggregate metrics had not 
previously distinguished. The deeper analysis of the collapse, however, 
qualifies its severity: for four of the five collapsed configurations 
the discriminative signal is present but masked by the $\arg\max$ 
(F10), and the mechanism responsible --- a response prior that dominates 
the signal --- both explains why monotone recalibration cannot help and 
predicts the re-thresholding remedy that does (F11). The discovery that 
one medical LLM (Meditron) ranks discriminative cases below chance (F9, 
F11) is more consequential than a mere accuracy failure: it is a 
corruption of the signal itself, distinct from the recoverable collapse 
of BioMistral, and it survives external replication on TREC (F12). The 
one silver lining among the pre-registered results --- the effectiveness 
of post-hoc isotonic recalibration (F6) --- applies only to models that 
classify; the complementary remedy for models that collapse is 
re-thresholding, not recalibration. The implications of these findings 
for clinical deployment and for future research are discussed in 
Chapter~\ref{chap:conclusion}.